% Shared preamble of the HSEMHS mode documents (docs/physics/*.tex, docs/technical/*.tex).
\documentclass[10pt]{extarticle}
\usepackage[a4paper,margin=1.7cm]{geometry}
\usepackage[T1]{fontenc}\usepackage{lmodern}\usepackage{tabularx}
\usepackage{amsmath,amssymb,bm,enumitem,titlesec,xcolor,booktabs,array,tikz,graphicx,longtable}
\usepackage[hidelinks]{hyperref}
\titleformat{\section}{\large\bfseries\color{blue!50!black}}{\thesection.}{0.5em}{}
\titlespacing*{\section}{0pt}{9pt}{3pt}
\titleformat{\subsection}{\normalsize\bfseries}{\thesubsection}{0.5em}{}
\titlespacing*{\subsection}{0pt}{6pt}{2pt}
\setlist{nosep,leftmargin=1.3em}
\setlength{\parindent}{0pt}\setlength{\parskip}{4pt}
\renewcommand{\arraystretch}{1.15}
\newcommand{\B}{\mathbf B}\newcommand{\J}{\mathbf J}\newcommand{\F}{\mathbf F}\newcommand{\R}{\mathbf R}
\newcommand{\zh}{\hat{\mathbf z}}\newcommand{\Bo}{\mathbf B_{\rm obs}}\newcommand{\Lv}{\mathbf L}
\newcommand{\code}[1]{\texttt{#1}}
\newcolumntype{L}[1]{>{\raggedright\arraybackslash}p{#1}}
\newcolumntype{Y}{>{\raggedright\arraybackslash}X}
\newcommand{\doctitle}[2]{{\Large\bfseries #1}\\[2pt]{\small #2}\par\vspace{2pt}}
\newif\ifadjointcolumn\newif\ifadjointpoisson   % which part of common/adjoint.tex a document shows
\pagestyle{plain}
\AtBeginDocument{\sloppy}

\begin{document}
\doctitle{column-hse-logm: hydrostatic columns on a column-mass grid}{Physics and method of \code{run.mode = "column-hse-logm"} in
HSEMHS, 25 Nov 2026. Built on the $\tau$-grid mode: read the \code{column-hse-ltau} document for the parts shared with it.}

\begin{abstract}
\noindent Some inversions return the temperature and the magnetic field on surfaces of constant column mass $m$ (the mass per unit
area above a point) instead of constant optical depth. This mode computes, for every pixel column, the gas pressure, density,
electron density, the 500\,nm extinction and optical depth, and the geometric height of every surface. On a mass axis hydrostatic
equilibrium is exact and needs no opacity: $P=gm$. The heights follow from the density, the optical depth becomes an output, and
the columns are placed against each other as in the $\tau$-grid mode.
\end{abstract}

\section{Introduction}
Write $m(z)=\int_z^\infty\rho\,dz'$ [kg\,m$^{-2}$], the mass above height $z$ per unit area. Its differential is
\begin{equation}
dm=-\rho\,dz ,\label{eq:dm}
\end{equation}
which has the same form as the optical depth, $d\tau=-\chi\,dz$: $\tau$ counts absorbers, $m$ counts mass. Everything the
\code{column-hse-ltau} document derives on the $\tau$ axis therefore holds on the mass axis with $\tau\to m$ and $\chi\to\rho$. The
one physical difference: the pressure of a column in hydrostatic equilibrium is simply the weight of the gas above,
\begin{equation}
\frac{dP}{dz}=-\rho g\quad\Longrightarrow\quad\frac{dP}{dm}=g\quad\Longrightarrow\quad P=g\,m ,\label{eq:pgm}
\end{equation}
with no reference to temperature, composition or opacity, provided no mass lies above the top level beyond its $m_1$ (the top
closure $P_1=gm_1$). \textbf{Input:} $\log_{10}m$ (common to all columns), $x$, $y$, $T$, $\B$. \textbf{Output:} $P$, $\rho$, $n_e$,
$\chi_{500}$, $\tau_{500}$ and the heights $z$ of every level of every column.

\section{Notation}
\begin{tabularx}{\linewidth}{l Y l}
\toprule
symbol & meaning & unit \\ \midrule
$m$ & column mass above a level (along the line of sight on input) & kg\,m$^{-2}$ \\
$\ell=\log_{10}m$ & the axis; levels $k=1\dots n$, top first & -- \\
$\rho$, $P$, $T$ & density, gas pressure, temperature & kg\,m$^{-3}$, Pa, K \\
$\chi=\chi_{500}$, $\kappa=\chi/\rho$ & extinction at 500\,nm, opacity per unit mass & m$^{-1}$, m$^2$\,kg$^{-1}$ \\
$g$, $\mu$ & gravity, cosine of the viewing angle & m\,s$^{-2}$, -- \\
\bottomrule
\end{tabularx}

% Shared section of the physics documents: equation of state and 500 nm extinction (\input by each mode document).
\section{The equation of state and the opacity}\label{sec:eos}
Every force balance in this document involves the density $\rho$, and every optical depth involves the extinction
coefficient $\chi_{500}$. Neither is an unknown of its own: at given temperature and pressure both are fixed by the
microphysics of the gas. This section explains how.

\subsection{The ideal gas}
In the solar photosphere and chromosphere the particles interact only weakly, so the gas obeys the ideal-gas law. If $n$
is the number of free particles (atoms, ions, molecules and electrons) per cubic metre,
\begin{equation}
P=n\,k_BT .\label{eq:ideal}
\end{equation}
If the average mass of a free particle is $\bar\mu\,m_{\rm u}$ (with $m_{\rm u}$ the atomic mass unit and $\bar\mu$ the
\textbf{mean molecular weight}), the mass density is $\rho=n\,\bar\mu\,m_{\rm u}$, and therefore
\begin{equation}
\rho=\frac{\bar\mu\,m_{\rm u}}{k_B}\,\frac{P}{T}.\label{eq:rho_ideal}
\end{equation}
At fixed $\bar\mu$ and $T$ the density is proportional to the pressure.

\subsection{\texorpdfstring{Why $\bar\mu$ depends on $T$ and $P$}{Why the mean molecular weight depends on T and P}: ionisation and molecules}
For neutral solar gas (about 91\,\% hydrogen and 9\,\% helium atoms by number, plus a small fraction of heavier elements)
$\bar\mu\approx1.3$. When a hydrogen atom is ionised it becomes two free particles, a proton and an electron, of almost
the same total mass; in fully ionised gas $\bar\mu\approx0.6$. In the coolest layers hydrogen partly forms H$_2$ molecules,
which raises $\bar\mu$. The degree of ionisation of an element in stage $i$ is given by the \textbf{Saha equation}
\begin{equation}
\frac{n_{i+1}\,n_e}{n_i}=\frac{2U_{i+1}(T)}{U_i(T)}\left(\frac{2\pi m_ek_BT}{h^2}\right)^{3/2}e^{-E_i/k_BT},\label{eq:saha}
\end{equation}
with $E_i$ the ionisation energy (13.6\,eV for hydrogen), $U_i$ the partition functions, $m_e$ the electron mass and $h$
Planck's constant. Two properties matter for us. (i) Because of the exponential, ionisation switches on over a narrow
range of temperature. (ii) The electron density $n_e$ appears on the left: at higher density there are more electrons to
recombine with, and a higher temperature is needed for the same ionisation. Hence $\bar\mu=\bar\mu(T,P)$, and
Eq.~\eqref{eq:rho_ideal} is no longer a simple proportionality. The Saha equations of all elements are coupled through the
common $n_e$ and through charge conservation ($n_e$ equals the total charge of the ions), and the molecular equilibria
(H$_2$, H$_2^+$, H$^-$) are solved together with them. This assumes \textbf{local thermodynamic equilibrium} (LTE): the
level populations and ionisation stages are those of a gas in equilibrium at the local $T$.

\subsection{\texorpdfstring{The 500\,nm extinction}{The 500 nm extinction}}
Light of wavelength 500\,nm travelling a distance $ds$ loses the fraction $\chi_{500}\,ds$ of its intensity, by absorption
and by scattering. In the solar photosphere $\chi_{500}$ is dominated by the negative hydrogen ion H$^-$ (a hydrogen atom
with a second, weakly bound electron, binding energy 0.75\,eV), whose abundance is proportional to the electron density.
The electrons come mostly from easily ionised metals (Na, Mg, Si, Ca, Fe, \dots), so $\chi_{500}$ grows steeply with
temperature and with density. The code adds the continuum contributions of H (bound--free and free--free), H$^-$,
H$_2^+$, He, He$^-$, the metals C, N, Mg, Si, Ca and Fe, Rayleigh scattering on H, He and H$_2$, and Thomson scattering
on free electrons, each counted once. Spectral lines are not included: $\chi_{500}$ is a continuum extinction.

\subsection{The table and its interpolation}
Solving the coupled Saha, molecular and charge equations at every grid point at every step would be far too slow. They
are therefore solved once, with Wittmann's equation-of-state code for solar composition (with a dedicated hot branch above
a switch temperature $T_{\rm sw}$, where the original solver fails in the hot, low-pressure corner), on a grid of
$X=\ln T$ and $Y=\ln P$, and stored in a table. For every node the table holds the three functions
\[
\ln\rho(X,Y),\qquad \ln n_e(X,Y),\qquad \ln\chi_{500}(X,Y),
\]
together with their derivatives $\partial_X$, $\partial_Y$ and $\partial_X\partial_Y$. Between the nodes each function is
evaluated by \textbf{bicubic Hermite interpolation}: a cubic polynomial in each direction that matches the stored values
and derivatives at the four corners of the cell. The result is continuous with continuous first derivatives, which is
important for the Newton-type solvers of the magnetohydrostatic modes, and it provides the derivative
$\partial\ln\rho/\partial\ln P|_T$ that they need. The table used by all modes is \code{cache/eos\_table\_fullz.h5}.
Its range of $T$ and $P$ (its \emph{envelope}) covers the input atmospheres; a point outside the envelope is reported as
an equation-of-state failure, never extrapolated.

Two derived quantities appear repeatedly:
\begin{equation}
q\equiv\frac{\rho}{P}\quad[{\rm s^2\,m^{-2}}],\qquad H\equiv\frac{P}{\rho\,g}=\frac{1}{g\,q}=\frac{k_BT}{\bar\mu\,m_{\rm u}\,g}\quad[{\rm m}].
\label{eq:H}
\end{equation}
$H$ is the \textbf{pressure scale height}. For $T=6000$\,K and $\bar\mu=1.3$, $H\approx140$\,km.


\section{The equations solved}
\textbf{Pressure.} Equation~\eqref{eq:pgm}. Along a ray at $\mu=\cos\theta<1$ the input axis is the mass along the ray,
$m_{\rm LOS}=m/\mu$, and $P=\mu g\,m_{\rm LOS}$. The code integrates the same two equations as on the $\tau$ axis,
\begin{equation}
\frac{d\ln P}{d\ln m}=\frac{\mu g\,m}{P},\qquad \frac{dz}{d\ln m}=-\frac{\mu\,m}{\rho},
\end{equation}
by Heun's method with the substeps of the $\tau$ mode; the first has the exact solution $P=\mu gm$ (its slope is 1 at every
substep), so only the heights need the equation of state: $\rho(T,P)$ from the table.

\textbf{Heights.} $z$ decreases strictly downward (Eq.~\ref{eq:dm}); one constant per column remains free, fixed by the alignment.

\textbf{Optical depth} is an output: $d\tau=\kappa\,dm$, i.e.\ $d\tau=(\chi/\rho)\,m\,d\ln m$, integrated level by level with the
trapezoidal rule (one interval per level, no substeps), from $\tau_1=\kappa_1m_1$ at the top (the closure of the $\tau$ axis:
$\kappa$ constant above the top level). Across a coarse level where the gas is hot above cool gas, this single trapezoid is less
accurate than the substepped integration of the $\tau$ mode.

\section{Where the columns sit}
As in \code{column-hse-ltau}: \code{align = "ptot"} shifts each column so that the total pressure $P+B^2/2\mu_0$ is equal at one
height (magnetic tension neglected); \code{align = "none"} puts $z=0$ at $\log_{10}m=0$ in every column. Then a global gauge sets the
mean height of $\tau_{500}=1$, computed from the output $\tau$, to zero. Absolute heights are relative to that gauge.

\section{The line of sight}
As on the $\tau$ axis: one ray per pixel on a common $\log_{10}m_{\rm LOS}$ axis; \emph{plane-parallel} rays are vertical columns
with $m_{\rm LOS}=m/\mu$; in the \emph{inclined} geometry the points of a ray are displaced by $z\tan\theta$, which does not change a
column-by-column balance, so both give the same result. $\mu=1$ is the vertical case exactly.

\section{Assumptions}
\begin{enumerate}[label=\textbf{A\arabic*.},leftmargin=2.2em]
\item \textbf{Static gas, constant vertical gravity, no mass above the top level beyond $m_1$.}
\item \textbf{Each column on its own.} No horizontal force balance; the alignment is a modelling choice.
\item \textbf{The mass axis is known.} An inversion that returns $m$ must get it right: $m$ fixes $P$ directly. A cgs axis
(g\,cm$^{-2}$) not marked as such scales every pressure by 10 and cannot be detected.
\item \textbf{LTE equation of state} (the NLTE table is an experimental option); $T$ and $\B$ fixed on the mass levels.
\end{enumerate}

\section{What the result means}
On the Bifrost snapshot (en024048, 385; input with $m$ from the snapshot's own $\rho$ and $P$, which an inversion does not have)
the heights of $\tau=1$ come out with an rms error of 56.6\,km against the snapshot, against 20.7\,km for \code{column-hse-ltau} on
the same columns; the reason is not yet understood. \code{mhs-poisson-logm} brings it to 19.5\,km.

\adjointcolumntrue
% Shared section of the physics documents of the six fixed-B modes: derivatives of the result (the adjoint).
% \input by each mode document after \adjointcolumntrue (column modes) or \adjointpoissontrue (Poisson modes).
\section{Derivatives of the result: the adjoint}\label{sec:adjoint}
An inversion that fits observed spectra changes $T$ (and $\B$) and asks how a misfit $\ell$ that depends on the atmosphere,
through $P$, $\rho$, $n_e$, $\chi_{500}$, $\tau_{500}$ and the heights $z$, changes with them. It needs the gradient
$\partial\ell/\partial\ln T$ at every grid point, often $10^6$ numbers. Perturbing each input in turn and solving again would
cost one or two solves per number. The \textbf{adjoint} gives all of them for the cost of about one solve.

\subsection{The idea}
Write the solve as a map from the inputs $x=(\ln T,\B)$ to the outputs $y$, and the misfit as $\ell(y)$. By the chain rule
\begin{equation}
\frac{\partial\ell}{\partial x}=\Big(\frac{\partial y}{\partial x}\Big)^{\!\top}\bar y,\qquad \bar y\equiv\frac{\partial\ell}{\partial y}.
\label{eq:adj}
\end{equation}
$\bar y$ (the \emph{cotangent} of the outputs) is given by the caller. The Jacobian $\partial y/\partial x$ is never formed:
the transposed product is evaluated by running the computation backwards, each step passing the cotangent of its outputs
to its inputs. This is reverse-mode differentiation, and its cost does not depend on how many inputs there are.

\subsection{What is differentiated}
\ifadjointcolumn
The column integration is a fixed sequence of arithmetic steps (substeps, table look-ups, the top closure, the alignment and
the gauge), so its adjoint is the derivative of exactly that computation: the \textbf{discrete adjoint}. Its discrete choices
are recorded during the solve and held fixed: the number of substeps of each interval, the table cell of each look-up, the
branch of the $\tau$ top closure, the alignment level and its median column, and the interpolation brackets. Away from the
points where one of these choices switches, the map is smooth and this is its exact derivative; at a switch no derivative
exists. Where a step length depends on the state (the Heun substeps of the $\tau$ and mass axes), the length is
differentiated too. An implicitly defined quantity (the hydrostatic top closure) is differentiated by the implicit function
theorem. The equation-of-state derivatives are those of the table interpolant (its bicubic Hermite polynomials), not of the model behind the table.
\fi
\ifadjointpoisson
The Poisson solve is iterative: its result is the state $v^*$ at which the gradient of $\Phi$ vanishes,
$G(v^*,x)\equiv\partial\Phi/\partial v=0$. Differentiating this condition (the \textbf{implicit function theorem}) gives
\begin{equation}
\frac{\partial v^*}{\partial x}=-\mathsf H^{-1}\,\frac{\partial G}{\partial x},\qquad
\mathsf H=\frac{\partial^2\Phi}{\partial v^2},\label{eq:ift}
\end{equation}
so the adjoint solves one linear system with the Hessian $\mathsf H$, $\mathsf H\lambda=\bar v$, and returns
$-(\partial G/\partial x)^{\top}\lambda$, plus the contributions of everything else that depends on $x$: the starting
column-hydrostatic state (which sets the weights of $\Phi$ through their layer medians), the part above the cut and the
gauge. The iterations themselves (damping, step scaling, preconditioners, warm starts) only steer the solver toward $v^*$
and are not differentiated. Two points matter:
\begin{itemize}
\item \textbf{The full Hessian.} With $\B$ fixed the residual does not vanish at the solution (the Lorentz force is in general
not a gradient), so the Gauss--Newton matrix $\mathsf J^{\top}\mathsf W\mathsf J$ of the solver is not the Hessian: the second
derivatives of the residuals contribute too. The adjoint uses the full $\mathsf H$, applied to vectors by differences of the
exact gradient, and solves for $\lambda$ by preconditioned conjugate gradients.
\item \textbf{Directions $\Phi$ does not see.} Some changes of the unknowns leave $\Phi$ unchanged and $\mathsf H$ singular along
them. They keep their starting values, so the derivative passes through to the start along them: on the height grid the mean
$\ln P$ of the anchor plane; on the $\tau$ and mass axes a uniform shift of all heights (fixed by the gauge) and the shortest
horizontal waves of the centred differences, which the differences cannot see.
\end{itemize}
Equation~\eqref{eq:ift} holds only at a stationary point. The Gauss--Newton iteration damps every step and measures its own
progress with that damping, so a slowly converging direction can stay far from stationary while the iteration reports
little left to gain. The option \code{polish} therefore adds up to $n$ \textbf{Newton steps with the exact Hessian} after the
solve, which bring the returned state to $v^*$ (on a small test cube four steps lower the undamped Newton decrement to
$10^{-25}$). The adjoint reports how stationary the returned state was: \code{:stationary} when one undamped Newton step
would lower $\Phi$ by at most a relative $10^{-6}$, \code{:not\_stationary} otherwise, in which case the result is the
derivative of the point the solver was approaching.
\fi

\subsection{How exact it is}
The adjoint is checked against finite differences of the solve itself: $\langle\bar y,\,[y(x+h\,\delta x)-y(x-h\,\delta x)]/2h\rangle$
against $\langle\partial\ell/\partial x,\,\delta x\rangle$ for random $\delta x$ and $\bar y$, on 1, 2 and 3 MPI ranks and every
viewing geometry.
\ifadjointcolumn
The two agree to $3\times10^{-10}$--$3\times10^{-8}$ (relative) on all three axes.
\fi
\ifadjointpoisson
With \code{polish} they agree to $10^{-9}$--$10^{-7}$ (relative).
\fi
A solve with a given electron density (no hydrostatic integration) is forward only and has no adjoint.


\section{Summary of the algorithm}
\begin{enumerate}
\item Read $\log_{10}m$, $x$, $y$, $T$, $\B$ (converting a cgs axis marked \code{units = "g cm-2"}).
\item Per column: $P=\mu g\,m_{\rm LOS}$; heights by Heun substeps of $dz/d\ln m=-\mu m/\rho$; $\rho$, $n_e$, $\chi$ from the table.
\item Align the columns (total pressure, or $z=0$ at $\log m=0$); $\tau_{500}$ by trapezoids of $\kappa\,dm$; gauge the mean
$z(\tau{=}1)$ to zero; check the fields; write.
\end{enumerate}
\end{document}
